\chapter{Dzienniczki}\markboth{Dzienniczki}{Dzienniczki}

Trzy tabele do wydruku. Można je drukować w tej postaci albo przenieść do arkusza, który
prowadzisz przez cały rok.

\section{Dziennik czytania}
Jeden wiersz na książkę, wpisany w dniu zakończenia, a nie w dniu, w którym ją zacząłeś. Ocena
w skali 1--5 ma sens tylko wtedy, gdy wiesz, co oznacza. Zapisz na marginesie jedną regułę, na
przykład: \glqq 5 = wróciłbym do tego tekstu\grqq{}, i trzymaj się jej.

{\renewcommand{\arraystretch}{2.1}\small
\begin{longtable}{@{}p{3.3cm}p{2.5cm}p{1.6cm}p{1.6cm}p{1.2cm}p{3.5cm}@{}}
\toprule
\textsf{\textbf{Tytuł}} & \textsf{\textbf{Autor}} & \textsf{\textbf{Start}} & \textsf{\textbf{Koniec}} & \textsf{\textbf{Ocena}} & \textsf{\textbf{Jedno zdanie}}\\
\midrule
\endhead
\bottomrule
\endfoot
& & & & & \\
& & & & & \\
& & & & & \\
& & & & & \\
& & & & & \\
& & & & & \\
& & & & & \\
& & & & & \\
& & & & & \\
& & & & & \\
& & & & & \\
& & & & & \\
& & & & & \\
& & & & & \\
& & & & & \\
& & & & & \\
\end{longtable}
}

\section{Siatka roku}
Pięćdziesiąt dwa tygodnie. W kolumnie \textbf{Książka} wpisz tytuł, który skończyłeś w tym
tygodniu, albo zostaw puste, jeśli nic nie skończyłeś. Kolumna \textbf{Minuty} to suma czasu
czytania w tygodniu; porównaj ją z wynikiem z dodatku B, zamiast pamiętać.

{\renewcommand{\arraystretch}{1.55}\small
\begin{longtable}{@{}p{2.0cm}p{5.6cm}p{2.2cm}p{2.2cm}p{1.6cm}@{}}
\toprule
\textsf{\textbf{Tydzień}} & \textsf{\textbf{Książka}} & \textsf{\textbf{Minuty}} & \textsf{\textbf{Strony}} & \textsf{\textbf{Gotowe}}\\
\midrule
\endhead
\bottomrule
\endfoot
1 & & & & \\
2 & & & & \\
3 & & & & \\
4 & & & & \\
5 & & & & \\
6 & & & & \\
7 & & & & \\
8 & & & & \\
9 & & & & \\
10 & & & & \\
11 & & & & \\
12 & & & & \\
13 & & & & \\
14 & & & & \\
15 & & & & \\
16 & & & & \\
17 & & & & \\
18 & & & & \\
19 & & & & \\
20 & & & & \\
21 & & & & \\
22 & & & & \\
23 & & & & \\
24 & & & & \\
25 & & & & \\
26 & & & & \\
27 & & & & \\
28 & & & & \\
29 & & & & \\
30 & & & & \\
31 & & & & \\
32 & & & & \\
33 & & & & \\
34 & & & & \\
35 & & & & \\
36 & & & & \\
37 & & & & \\
38 & & & & \\
39 & & & & \\
40 & & & & \\
41 & & & & \\
42 & & & & \\
43 & & & & \\
44 & & & & \\
45 & & & & \\
46 & & & & \\
47 & & & & \\
48 & & & & \\
49 & & & & \\
50 & & & & \\
51 & & & & \\
52 & & & & \\
\end{longtable}
}

\section{Lista do przeczytania}
Tytuły, które chcesz przeczytać. Kolumna \textbf{Próbka} to liczba stron, które przeczytasz,
zanim podejmiesz decyzję. Decyzja to jedno z trzech słów: \emph{czytam}, \emph{odkładam} albo
\emph{porzucam}. Odłożenie nie jest porażką, ale tytuł odłożony bez daty wraca na listę za rok.

{\renewcommand{\arraystretch}{2.0}\small
\begin{longtable}{@{}p{3.4cm}p{2.5cm}p{3.2cm}p{2.0cm}p{2.8cm}@{}}
\toprule
\textsf{\textbf{Tytuł}} & \textsf{\textbf{Autor}} & \textsf{\textbf{Skąd o niej wiesz}} & \textsf{\textbf{Próbka (strony)}} & \textsf{\textbf{Decyzja}}\\
\midrule
\endhead
\bottomrule
\endfoot
& & & & \\
& & & & \\
& & & & \\
& & & & \\
& & & & \\
& & & & \\
& & & & \\
& & & & \\
& & & & \\
& & & & \\
& & & & \\
& & & & \\
\end{longtable}
}
